\section{Sistema Integral de Monitoreo Ambiental (SIMA)}
\label{sec:sima}

El Sistema Integral de Monitoreo Ambiental (SIMA) es la red del Gobierno de
Nuevo León encargada de medir y vigilar la calidad del aire. Su objetivo es
conocer las concentraciones de contaminantes a las que está expuesta la
población, generar alertas cuando existen condiciones de contaminación elevada
y mantener informados a los ciudadanos sobre la salud ambiental y la calidad del
aire \parencite{gobiernoNL2026PIGECA}.

SIMA inició operaciones el 20 de noviembre de 1992 como un departamento de la
Subsecretaría de Protección al Medio Ambiente y Recursos Naturales, con tan
solo cinco estaciones de monitoreo. Durante 32 años ha crecido y modernizado
su infraestructura \parencite{villasaez2024}. Actualmente cuenta con 15
estaciones fijas distribuidas estratégicamente en 11 municipios de la zona
metropolitana de Monterrey, además de 2 estaciones móviles
\parencite{gobiernoNL2026ModernizacionSIMA}.

SIMA mide ocho contaminantes criterio y siete parámetros meteorológicos,
detallados junto con sus unidades y cobertura en el diccionario de datos
(\cref{tab:dic-contaminantes,tab:dic-meteo}).

SIMA es la infraestructura gubernamental encargada de convertir la información
sobre la contaminación atmosférica en acciones de gestión ambiental, y es el
medio por el cual se obtiene una parte fundamental de la información necesaria
para hacerlo.

Los datos utilizados en este trabajo se relacionan directamente con SIMA, ya
que provienen de las mediciones realizadas por sus estaciones de monitoreo
ambiental. Las estaciones registran regularmente contaminantes y variables
meteorológicas. Una vez extraídos, los datos se procesan y validan para
integrarlos en las bases de datos de la institución.